\section{Does the Model Depend on Lung Content?}
\label{sec:lung-dependence}

On the patient-disjoint, deduplicated split the vision model reaches AUROC 0.9978, yet only about half of the
Grad-CAM activation in nonzero maps falls inside the segmented lungs (Section~\ref{sec:dedup}). This section
asks (i) whether classification depends on lung content or can be carried by non-lung regions and acquisition
cues, (ii) whether the Grad-CAM lung fraction tracks that dependence, and (iii) whether lung-aware training
changes it. All analyses use the same 1,165 test images (849 Pneumonia). Unless stated otherwise they use the
single frozen vision checkpoint (seed 42). Intervals are image-level bootstrap intervals (2,000 resamples)
and ignore the clustering of images within patients. Interventions replace pixels with mid-gray, which the
model did not see in training and which pushes its outputs toward Pneumonia; we therefore compare AUROC
rather than sensitivity or specificity, and pair each lung intervention with a displaced-mask control of the
same shape and area (mean overlap with the true mask 0.26). The analyses are descriptive: no multiplicity
adjustment was applied, and none of them establishes clinical validity.

\subsection{Lung masks}
Masks come from an automatic segmenter trained on adult radiographs (mean area 0.38 of the image; 1,162 of 1,165
masks cover between 10\% and 70\% of it). In a visual check of 24 random test images the outline followed
aerated lungs well, but 6--8 tiles showed defects, clustered on consolidated lungs: the outline cut into a
dense lower zone, left a notch, or dropped a region. The mask also excludes the mediastinum and hilum, so
``background'' contains the heart, ribs, shoulders, abdomen and text markers. We therefore use the masks for
interventions and report a dilation sensitivity check; we do not treat them as expert annotations.

\subsection{Frozen-model interventions}
Table~\ref{tab:audit} gives the audit. Keeping only the lungs gives AUROC 0.9746, compared with 0.9374 when the
same-shaped displaced region is kept; graying the lungs costs 0.022, about four times the 0.005 lost when the
displaced region is grayed. Lung content is therefore informative. It is not necessary: with the lungs grayed
the model still reaches 0.9759, and graying the outer 10\% margin changes AUROC by $-0.0007$ (interval
includes zero). Growing the masks by 4 and 8 px (Euclidean) moved keep-lungs AUROC to 0.982 and 0.989 and
gray-lungs AUROC to 0.970 and 0.964. With the 8~px masks covering 52\% of the image, the model still reaches
0.964 with that region grayed, so segmentation defects at the lung boundary do not explain the non-lung signal.

\begin{table}[htbp]
\centering
\caption{Frozen-model interventions on the test set ($n=1{,}165$). Matched controls use a displaced copy of
the lung mask (mean of five random shifts per image).}
\label{tab:audit}
\small
\begin{tabular}{lcc}
\toprule
Condition & AUROC [95\% CI] & Change vs original [95\% CI] \\
\midrule
Original image & 0.9978 [0.9962, 0.9990] & --- \\
Keep lungs only & 0.9746 [0.9650, 0.9829] & $-0.0232$ [$-0.0324$, $-0.0155$] \\
\quad matched: keep displaced region only & 0.9374 [0.9229, 0.9509] & $-0.0604$ [$-0.0748$, $-0.0468$] \\
Gray the lungs & 0.9759 [0.9672, 0.9839] & $-0.0219$ [$-0.0304$, $-0.0144$] \\
\quad matched: gray displaced region & 0.9929 [0.9892, 0.9959] & $-0.0049$ [$-0.0079$, $-0.0025$] \\
Gray the outer 10\% margin & 0.9971 [0.9949, 0.9988] & $-0.0007$ [$-0.0021$, $+0.0005$] \\
\bottomrule
\end{tabular}
\end{table}

\subsection{Acquisition cues}
The two classes were exported differently. In the test split the median Normal file is $1644\times1314$ px and
495~KB, the median Pneumonia file $1160\times776$ px and 71~KB, and no image size shared by at least 15 test
images contains both classes. Classifiers trained on the same training split but given only such properties
recover the label almost as well as the network (Table~\ref{tab:cue-baselines}); file size alone reaches
0.980. This shows that the label is recoverable from acquisition properties, a precondition for shortcut
learning, not that the network uses them: it receives $224\times224$ resized, normalized pixels and cannot read
file size or original dimensions. File size is also not a pure artefact, since compressibility may carry some
disease signal.

\begin{table}[htbp]
\centering
\caption{Test AUROC of classifiers that see only acquisition-level features (model chosen on validation data).}
\label{tab:cue-baselines}
\small
\begin{tabular}{lc}
\toprule
Features & Test AUROC [95\% CI] \\
\midrule
File size (KB) alone & 0.980 [0.972, 0.988] \\
Width, height, aspect ratio & 0.916 [0.901, 0.932] \\
Mean and SD of pixel intensity & 0.749 [0.718, 0.778] \\
$16\times16$ gray thumbnail (256 values) & 0.977 [0.968, 0.985] \\
Size, file size and intensity statistics (18 features) & 0.986 [0.979, 0.992] \\
DenseNet-121 + CBAM (vision only) & 0.998 \\
\bottomrule
\end{tabular}
\end{table}

The network still outperforms every baseline (smallest paired gap 0.012, bootstrap interval [0.006, 0.019]).
On the 55 test images where the file-size classifier errs (at its validation Youden threshold), the
network's AUROC falls to 0.923 [0.845, 0.978] and its error rate is 20.0\%, against 1.1\% on the remaining
1,110; 11 of its 23 test errors lie among those 55. This is correlational: atypical images may be hard for both.

To test the network's sensitivity to such cues directly we perturbed the original images before the usual
resize (Table~\ref{tab:cue-perturbation}). JPEG recompression (quality 75 and 25) and a $2\times$ down-up
resampling left AUROC unchanged, so compression artefacts and fine texture do not drive the ranking. Histogram
equalization lowered AUROC by 0.007 overall and by 0.113 [$-0.235$, $-0.011$] on the 55 discordant images,
the only result pointing toward reliance on global intensity; it also produces intensity distributions the
model never saw, so it is not a clean removal of one cue. Padding to a square kept AUROC at 0.991 but raised
the false-negative rate at the 0.5 threshold from 1.4\% to 29.7\% (mean logit shift $-5.1$): framing shifts
the calibration of all images without changing the ranking much.

\begin{table}[htbp]
\centering
\caption{Perturbing acquisition cues in the original image (frozen checkpoint, $n=1{,}165$).}
\label{tab:cue-perturbation}
\small
\begin{tabular}{lccc}
\toprule
Perturbation & AUROC [95\% CI] & Change [95\% CI] & FNR at 0.5 \\
\midrule
None & 0.9978 [0.9963, 0.9990] & --- & 1.4\% \\
JPEG quality 25 & 0.9978 [0.9962, 0.9990] & $0.0000$ [$-0.0003$, $+0.0003$] & 1.2\% \\
Histogram equalization & 0.9905 [0.9861, 0.9943] & $-0.0073$ [$-0.0113$, $-0.0038$] & 5.5\% \\
Pad to square & 0.9907 [0.9864, 0.9943] & $-0.0071$ [$-0.0107$, $-0.0039$] & 29.7\% \\
Down-up $2\times$ & 0.9978 [0.9963, 0.9990] & $0.0000$ [$-0.0002$, $+0.0002$] & 1.3\% \\
\bottomrule
\end{tabular}
\end{table}

\subsection{Does the Grad-CAM lung fraction track the interventions?}
Among the 847 Pneumonia images with a nonzero map, the lung fraction of Grad-CAM mass is strongly tied to
mask area (Spearman 0.83), so all correlations below are partial Spearman coefficients controlling for mask
area, with image-level bootstrap intervals. Grad-CAM mass is enriched inside the lungs: the median ratio of
lung fraction to mask area is 1.26, and 97\% of maps exceed 1. Images whose Grad-CAM mass lies more inside the
lungs lose more of their Pneumonia margin when the lungs are grayed ($\rho=-0.634$ [$-0.677$, $-0.582$]), while
the logit with only the lungs kept is unrelated to the fraction ($\rho=-0.025$ [$-0.095$, $+0.046$]). Part of
this is generic: the matched displaced-mask contrast correlates too ($+0.333$ [$+0.266$, $+0.393$]), and the
original logit itself correlates with the fraction ($-0.415$). The part specific to the true lung mask is the
difference between the two contrasts, $+0.203$ [$+0.129$, $+0.281$]. The lung fraction therefore reflects how
much of the prediction survives without the lungs, but only partly, and the analysis is correlational. Normal
images are not analysed because only 50 have a nonzero map.

\subsection{Lung-aware training}
\label{sec:lung-aware}
To test whether reliance on non-lung regions can be reduced, we retrained the same DenseNet-121 + CBAM model
(same split, hyperparameters and early-stopping rule) on masked images, with three seeds (42, 123, 2024) per
variant. The mask, grown by 8~px, is applied after resizing and before augmentation, with mid-gray outside.
\emph{Thorax-only} masks every training and validation image; \emph{random-background} masks each training
image with probability 0.5 and validates on full images; \emph{fixed-box} is an outline control that shows every
image through the same centred box (72\% of each side, area 0.52, matching the mean thorax mask), so its
cut-out carries no per-image information. Best validation AUROC was 0.9965--0.9968 (baseline), 0.9958--0.9966
(random-background), 0.9953--0.9966 (thorax-only) and 0.9953--0.9967 (fixed-box), so validation does not
separate them. Table~\ref{tab:lung-aware} reports test AUROC when only the named region is kept.

\begin{table}[htbp]
\centering
\caption{Test AUROC when only the named region is kept (rest mid-gray), mean $\pm$ SD over three seeds.
The baseline column differs slightly from the frozen checkpoint (0.9978) because it averages seeds 42, 123 and
2024 (0.9978, 0.9961, 0.9982). Dashes: not evaluated.}
\label{tab:lung-aware}
\footnotesize
\setlength{\tabcolsep}{4pt}
\begin{tabular}{lcccc}
\toprule
Input shown & Baseline & Random-bg. & Thorax-only & Fixed-box \\
\midrule
Original image & $0.9973\pm0.0011$ & $0.9974\pm0.0005$ & $0.9733\pm0.0157$ & $0.9937\pm0.0006$ \\
Dilated thorax only & $0.9879\pm0.0030$ & $0.9972\pm0.0004$ & $0.9975\pm0.0009$ & $0.9949\pm0.0014$ \\
Fixed central box only & $0.9933\pm0.0019$ & $0.9956\pm0.0016$ & $0.9944\pm0.0019$ & $0.9965\pm0.0012$ \\
Outside thorax only & $0.9672\pm0.0033$ & $0.9572\pm0.0053$ & $0.7441\pm0.0959$ & --- \\
Outside thorax, outer margin removed & $0.8479\pm0.0286$ & $0.8893\pm0.0195$ & $0.8628\pm0.0343$ & --- \\
Outer 10\% margin only & $0.9419\pm0.0069$ & $0.9355\pm0.0208$ & $0.6789\pm0.0722$ & --- \\
Outside the fixed box & $0.9578\pm0.0070$ & $0.9445\pm0.0245$ & $0.7132\pm0.0550$ & $0.8107\pm0.0347$ \\
\bottomrule
\end{tabular}
\end{table}

Three observations follow. First, a model that sees only the dilated thorax matches the full-image model
(0.9975 against 0.9973). Second, random-background training makes the model read the thorax as well as the
full image (0.9972 against 0.9879 for the baseline) but does not remove the outside signal (0.957 against
0.967 for the region outside the thorax; 0.936 against 0.942 for the margin alone), differences of about 0.01
or less that are comparable to seed-to-seed variation; training on full images half the time leaves the other cues available. Third, thorax-only training
removes most of the outside signal (0.744, seed range 0.685--0.855), but the size of this drop is uncertain.
The thorax-only models are out of distribution on unmasked images, which explains their lower and more
variable original-image AUROC.

A thorax-only model sees the lungs through a gray cut-out, whose outline could itself carry the label. A
classifier given only the shape of the lung mask (area, bounding box, row and column profiles, or a
$16\times16$ thumbnail of the mask, no pixel intensities) reaches test AUROC 0.946 [0.932, 0.959] (0.938 for
the 8~px masks used in training), although mask area alone gives only 0.61. Aspect ratio explains part of this,
because images are resized to a square: aspect ratio alone gives 0.844, and the outline classifier still
reaches 0.81--0.96 inside aspect-ratio quintiles (the two highest-aspect bins contain only 16 and 6 Normal
images, so those values are noisy). The fixed-box control shows that the outline is not needed for the
network's AUROC: models that see the identical box for every image reach 0.9965 on that input, level with the
thorax-only models on theirs, and the baseline already scores 0.9933 on the box. The residual signal outside
the thorax is also not an outline effect, since the fixed-box models reach 0.81 on the region outside their box
without any per-image outline. The thorax-only and fixed-box models differ by 0.001, inside seed spread, so
this experiment cannot say whether the thorax-only models use the outline at all.

\subsection{Summary}
Test AUROC of about 0.997 is reproduced from the dilated thorax alone, from a fixed central box alone, and (at
0.96) from the region outside the lungs alone, while the label is recoverable from file size alone (0.980).
The network is not insensitive to lung content, since removing the lungs costs four times what removing a
matched displaced region costs and the lung fraction of Grad-CAM tracks part of that effect, but no single
region, outline or cue is necessary. The 0.998 AUROC therefore says little about which evidence the model
uses. The experiments neither show that any of these models reads lung pathology nor support a claim of
clinical validity, and random-background training did not reduce reliance on non-lung regions, whereas
thorax-only training removed most of it by construction.
